\scanpage{29}
on s'intéressera plus particulièrement, pour des entiers $n$ convenables, aux repr. de $D_\infty$ qui veulent bien se factoriser par $D_n$. C'est dire que, \deleted{\unsure{le}} de ce point de vue, la question de fidélité ne se pose plus -- les représentations les plus intéressantes \inserted{\text{pour nous}} de $D_\infty$ \deleted{étant} -- celles qui donnent naissance aux \deleted{\unsure{correspondent}} polygônes réguliers finis -- étant précisément non fidèles!

% slanted
\margin{\uline{Pt de vue} $\Pi$-\uline{épinglé} $\Pi$ polygone comb. quelconque (fini ou infini)}
Arrivé à ce point, il y a lieu de considérer -- de $=$ que des polygônes \uline{combinatoires} infinis \inserted{\text{$\Pi=(S,\omega)$ --}} (tout comme les poly comb. finis) des réalisations \inserted{\text{ou \uline{repr.}}} géom de ceux-ci (pas néc. fidèles), et de $=$ des repr. géométriques \uline{régulières}, \deleted{en exigeant pour ces dernières} \deleted{\unsure{Une réalisation}} \deleted{en exigeant \unsure{pour} ces dernières} \deleted{que} % slanted
\margin{\uline{Déf}} On entendra par là une
% archive sheet number: 28
